% Компиляция: XeLaTeX. Кодировка: UTF-8.
\documentclass[a4paper,14pt]{extarticle}

\usepackage{fontspec}
\setmainfont{Times New Roman}
\usepackage[russian]{babel}
\addto\captionsrussian{\renewcommand{\contentsname}{Содержание}}
\usepackage[left=25mm,right=20mm,top=25mm,bottom=20mm]{geometry}
\usepackage{setspace}
\usepackage{titlesec}
\usepackage[titles]{tocloft}
\renewcommand{\cftsecleader}{\cftdotfill{\cftdotsep}}
\usepackage{array}

\onehalfspacing
\setlength{\parindent}{1.25cm}
\setcounter{secnumdepth}{0}
\titleformat{\section}{\normalfont\normalsize\bfseries}{}{0pt}{}
\titlespacing*{\section}{\parindent}{1.5\baselineskip}{0.5\baselineskip}

% Сведения для титульного листа ВКР. Неизвестные данные заполните здесь.
\newcommand{\StudentName}{Н.Ф. Михайлов}
\newcommand{\DegreeName}{БАКАЛАВРА} % При необходимости: МАГИСТРА или СПЕЦИАЛИСТА.
\newcommand{\DirectionName}{09.03.02 --- Информационные системы и технологии}
\newcommand{\ProfileName}{Информационно-управляющие системы}
\newcommand{\FacultyName}{ФКТИ}
\newcommand{\DepartmentName}{ИСиТ}
\newcommand{\DepartmentHead}{Е.С. Новикова}
\newcommand{\ThesisTitle}{Разработка системы автоматизированного нормоконтроля академических документов}
\newcommand{\SupervisorName}{И.О. Фамилия}
\newcommand{\ConsultantOne}{Н.В. Дюков}
\newcommand{\ConsultantTwo}{И.О. Фамилия}
\newcommand{\ConsultantThree}{И.О. Фамилия}

% Подпись и пояснения набираются внутри ячейки, чтобы не смещать столбцы.
\newcommand{\TitleSignature}{%
    \parbox[t]{\linewidth}{\centering
        \rule{\linewidth}{0.4pt}\\[-3pt]
        {\fontsize{8}{10}\selectfont\itshape подпись}}}
\newcommand{\TitleDegree}{%
    \parbox[t]{\linewidth}{\centering
        \vspace{14pt}
        {\fontsize{8}{10}\selectfont\itshape (Уч. степень, уч. звание)}}}

\begin{document}

% Поля образца Word: слева 30 мм, справа 15 мм, сверху и снизу 20 мм.
\newgeometry{left=30mm,right=15mm,top=20mm,bottom=20mm}
\begin{titlepage}
    \setstretch{1.4}
    \setlength{\parindent}{0pt}
    \centering
    {\fontsize{13}{15.6}\selectfont\bfseries
        Санкт-Петербургский государственный электротехнический университет\\
        «ЛЭТИ» им. В.И. Ульянова (Ленина)\\
        (СПбГЭТУ «ЛЭТИ»)\par}

    \vspace{7mm}

    {\setlength{\tabcolsep}{0pt}
    \renewcommand{\arraystretch}{1.05}
    \begin{tabular}{@{}>{\raggedright\arraybackslash\bfseries}p{0.435\textwidth}%
                        >{\raggedright\arraybackslash}p{0.565\textwidth}@{}}
        Направление (специальность) & \DirectionName \\
        Профиль (программа, специализация) & \ProfileName \\
        Факультет/Институт & \FacultyName \\
        Кафедра & \DepartmentName
    \end{tabular}\par}

    \vspace{7mm}
    {\setlength{\tabcolsep}{0pt}
    \begin{tabular}{@{}p{0.435\textwidth}p{0.228\textwidth}%
                        >{\raggedleft\arraybackslash}p{0.337\textwidth}@{}}
        Зав. кафедрой & \rule{0.95\linewidth}{0.4pt} & \DepartmentHead
    \end{tabular}\par}

    \vspace{17mm}

    {\bfseries ВЫПУСКНАЯ КВАЛИФИКАЦИОННАЯ РАБОТА\\
        \DegreeName\par}

    \vspace{7mm}
    {Тема: \ThesisTitle\par}

    \vspace{17mm}

    {\singlespacing
    \setlength{\tabcolsep}{0pt}
    \begin{tabular}{@{}p{0.217\textwidth}p{0.22\textwidth}p{0.264\textwidth}%
                        p{0.028\textwidth}>{\raggedright\arraybackslash}p{0.271\textwidth}@{}}
        Студент & & \TitleSignature & & \StudentName \\[5mm]
        Руководитель & \TitleDegree & \TitleSignature & & \SupervisorName \\[5mm]
        Консультанты & \TitleDegree & \TitleSignature & & \ConsultantOne \\[5mm]
        & \TitleDegree & \TitleSignature & & \ConsultantTwo \\[5mm]
        & \TitleDegree & \TitleSignature & & \ConsultantThree
    \end{tabular}\par}

    \vfill
    Санкт-Петербург\\
    2026
\end{titlepage}
\restoregeometry

\setcounter{page}{2}
\tableofcontents
\clearpage

\section{Введение}
В процессе подготовки выпускных квалификационных работ, курсовых проектов,
отчётов по практике и других академических документов обучающиеся должны соблюдать установленные требования к их структуре и оформлению.
Данные требования определяются ГОСТами, а также локальными нормативными документами и методическими указаниями образовательных организаций.

Нормоконтроль представляет собой процесс проверки соответствия документа установленным требованиям.
Он включает контроль параметров форматирования текста, оформления заголовков, таблиц, рисунков, списков, библиографических записей
и других структурных элементов документа.

Традиционно нормоконтроль осуществляется путём ручной проверки документа преподавателем, нормоконтролёром или самим обучающимся.
Однако развитие информационных технологий позволяет автоматизировать значительную часть данного процесса.
В настоящее время существуют программные решения, обеспечивающие автоматическую проверку оформления академических документов,
выявление нарушений и формирование отчётов с замечаниями.

При этом требования к оформлению академических документов могут различаться в зависимости от образовательной организации,
факультета, кафедры и типа работы, что необходимо учитывать при разработке и использовании автоматизированных систем нормоконтроля.
\section{Актуальность}
Актуальность работы обусловлена высокой трудоёмкостью ручного нормоконтроля академических документов,
вероятностью пропуска нарушений оформления, а также необходимостью учитывать различия в требованиях образовательных организаций.

Несмотря на существование программных решений для автоматизации нормоконтроля, возможности их адаптации
к индивидуальным требованиям пользователей различаются.
Не все рассмотренные системы предоставляют инструменты для самостоятельного изменения параметров проверки,
управления набором правил и создания пользовательских профилей требований.

В связи с этим актуальной является разработка системы автоматизированного нормоконтроля академических документов,
обеспечивающей гибкую настройку критериев проверки, выявление нарушений оформления и формирование понятных отчётов.
Использование такой системы позволит сократить трудозатраты на проверку документов, сделать процесс более прозрачным
и обеспечить применение различных наборов требований без необходимости изменения программного кода.
\section{Анализ конкуретнов}
Анализ существующих программных решений показал, что рынок автоматизированного нормоконтроля академических документов представлен сервисами,
ориентированными как на индивидуальных пользователей, так и на образовательные организации.
Рассмотренные решения обеспечивают автоматическую проверку параметров форматирования, структуры документа и оформления библиографических элементов.
Некоторые системы также поддерживают локальные требования кафедр и университетов.
Вместе с тем по открытой документации не для всех решений подтверждена возможность самостоятельного и детального управления набором правил проверки.
Это позволяет обосновать разработку NormLint как системы с конфигурируемыми профилями нормоконтроля, прозрачными результатами анализа
и возможностью быстрой и удобной адаптации к требованиям различных образовательных организаций.
\end{document}
